\bmtchapitre{Organisation des données et statistique}{Calculer proportions et évolutions, résumer une série par ses indicateurs et interpréter une boîte à moustaches.}{Données et probabilités}
\label{t11s-c18}
\begin{revision}Pourcentages, effectifs, fréquences, ordre des valeurs et moyenne pondérée.\end{revision}
\begin{automatismes}\exo\label{t11s-c18-auto}Calculer $20\%$ de $150$; ranger $8,3,5,3$; calculer leur moyenne.\end{automatismes}
\begin{corrige}[exercice~\ref{t11s-c18-auto}]\label{sol-t11s-c18-auto}$30$; $3,3,5,8$; moyenne $19/4=4{,}75$.\end{corrige}

\section{Proportion et évolution}
\begin{definition}Une proportion est un rapport partie/ensemble, avec effectif total non nul. Le taux d’évolution d’une valeur initiale $V_i>0$ à une valeur finale $V_f$ est $t=(V_f-V_i)/V_i$. Le coefficient multiplicateur est $1+t$; un taux de $p\%$ vaut $p/100$.\end{definition}
\begin{propriete}Deux évolutions successives de taux $t_1,t_2$ ont pour coefficient $(1+t_1)(1+t_2)$; leur taux global est ce produit moins $1$. Si $1+t\ne0$, le taux réciproque est $1/(1+t)-1$.\end{propriete}
\begin{demonstration}La valeur finale après deux évolutions vaut $V_i(1+t_1)(1+t_2)$. Pour revenir en arrière, multiplier par l’inverse du coefficient.\end{demonstration}
\begin{exemple}Hausse de $20\%$ puis baisse de $20\%$ : coefficient $1{,}2\times0{,}8=0{,}96$, baisse globale $4\%$. Passer de $30\%$ à $40\%$ est une hausse de $10$ points de pourcentage, mais une hausse relative de $1/3$, soit environ $33{,}3\%$.\end{exemple}
\section{Indicateurs de position}
\begin{definition}Pour une série ordonnée $x_1\le\cdots\le x_N$, $N\ge1$, la moyenne est $\bar x=\frac1N\sum x_i$; avec des effectifs $n_j$, elle vaut $\sum n_jx_j/\sum n_j$. La médiane est la valeur centrale si $N$ est impair, la moyenne des deux centrales si $N$ est pair.\end{definition}
\begin{definition}[convention de rang]
$Q_1$ est la valeur de rang $\lceil N/4\rceil$, $Q_3$ celle de rang $\lceil3N/4\rceil$. Le décile $D_k$ est la valeur de rang $\lceil kN/10\rceil$, $k=1,\ldots,9$. $\lceil a\rceil$ désigne le plus petit entier supérieur ou égal à $a$. Ces conventions sont précisées car elles peuvent différer selon les logiciels.
\end{definition}
\section{Dispersion et représentation}
\begin{definition}L’étendue vaut maximum moins minimum; l’écart interquartile vaut $Q_3-Q_1$. La variance descriptive vaut $V=\frac1N\sum(x_i-\bar x)^2$; l’écart type est $\sigma=\sqrt V$. Une boîte à moustaches représente minimum, $Q_1$, médiane, $Q_3$, maximum, selon la convention utilisée ici.\end{definition}
\begin{propriete}$V=\frac1N\sum x_i^2-\bar x^2$. La variance est positive ou nulle, et vaut zéro exactement lorsque toutes les valeurs sont égales.\end{propriete}
\begin{demonstration}Développer les carrés et utiliser $\sum x_i=N\bar x$. La première définition est une moyenne de carrés positifs ou nuls; elle est nulle exactement si chaque carré est nul.\end{demonstration}
\begin{exemple}[huit temps en minutes]
Série $4,6,6,8,10,10,12,14$. $N=8$, somme $70$, moyenne $8{,}75$, médiane $9$; rangs $2$ et $6$, donc $Q_1=6,Q_3=10$. Étendue $10$, écart interquartile $4$. Somme des carrés $792$; variance $99-(8{,}75)^2=22{,}4375$; écart type $\sqrt{22{,}4375}\simeq4{,}737$ minutes. La variance s’exprime en minutes carrées.
\end{exemple}
\begin{visualisation}[la boîte ne montre pas chaque donnée]
\begin{center}\begin{tikzpicture}[x=0.55cm,y=1cm]
\draw[->](3,0)--(15,0)node[right]{minutes};
\foreach \x in {4,6,8,10,12,14}{\draw(\x,0.06)--(\x,-0.06)node[below]{$\x$};}
\draw[bmtPalissandre,thick](4,0.9)--(6,0.9);\draw[bmtPalissandre,thick](10,0.9)--(14,0.9);
\draw[bmtPalissandre,thick](6,0.55)rectangle(10,1.25);\draw[bmtOcre,thick](9,0.55)--(9,1.25);
\draw(4,0.7)--(4,1.1);\draw(14,0.7)--(14,1.1);
\end{tikzpicture}\end{center}
Le trait intérieur est la médiane $9$, et non la moyenne $8{,}75$.
\end{visualisation}
\section{Choisir un résumé sans surinterpréter}
Une moyenne sensible aux valeurs extrêmes et une médiane n’apportent pas la même information. Comparer des séries avec mêmes unités et conventions. Pour des données groupées en classes, remplacer chaque valeur par le centre de sa classe donne une moyenne et une variance approchées; les données originales ne peuvent être reconstruites exactement.

\section{Exercices et problèmes}
\exo[Taux successifs]\label{t11s-c18-e01}
Un prix de $50000$ ariary augmente de $10\%$, puis diminue de $5\%$. Calculer le prix final et le taux global.
\begin{corrige}[exercice~\ref{t11s-c18-e01}]\label{sol-t11s-c18-e01}
$50000\times1{,}1\times0{,}95=52250$ ariary. Coefficient $1{,}045$, hausse globale $4{,}5\%$.
\end{corrige}
\exo[Retour au départ]\label{t11s-c18-e02}
Après une hausse de $25\%$, quelle baisse ramène à la valeur initiale ?
\begin{corrige}[exercice~\ref{t11s-c18-e02}]\label{sol-t11s-c18-e02}
Coefficient retour $1/1{,}25=0{,}8$, donc baisse $20\%$. Les pourcentages portent sur des bases différentes.
\end{corrige}
\exo[Effectifs]\label{t11s-c18-e03}
Les valeurs $5,10,15$ ont respectivement les effectifs $2,5,3$. Calculer moyenne, médiane, quartiles, $D_1,D_9$, variance et écart type.
\begin{corrige}[exercice~\ref{t11s-c18-e03}]\label{sol-t11s-c18-e03}
$N=10$, somme $105$, moyenne $10{,}5$. Série ordonnée : deux $5$, cinq $10$, trois $15$; médiane $10$, $Q_1=10$ (rang $3$), $Q_3=15$ (rang $8$), $D_1=5$, $D_9=15$. Moyenne des carrés $(50+500+675)/10=122{,}5$; variance $12{,}25$, écart type $3{,}5$.
\end{corrige}
\exo[Mêmes moyennes]\label{t11s-c18-e04}
Comparer les séries $(8,10,12)$ et $(0,10,20)$ : moyenne, médiane, étendue et variance. Que montre la comparaison ?
\begin{corrige}[exercice~\ref{t11s-c18-e04}]\label{sol-t11s-c18-e04}
Moyenne et médiane $10$ pour les deux. Étendues $4$ et $20$; variances $8/3$ et $200/3$. Une même moyenne ne garantit pas une même dispersion.
\end{corrige}
\exo[Décalage]\label{t11s-c18-e05}
Si on ajoute $3$ à toutes les valeurs d’une série, que deviennent moyenne, médiane, variance et écart type ?
\begin{corrige}[exercice~\ref{t11s-c18-e05}]\label{sol-t11s-c18-e05}
Moyenne et médiane augmentent de $3$. Les écarts à la moyenne restent identiques : variance et écart type inchangés. L’ordre des données est conservé.
\end{corrige}
\exo[Boîte à moustaches]\label{t11s-c18-e06}
Pour la série $2,3,4,4,5,7,8,9,10,12$, calculer les cinq valeurs nécessaires à la boîte et la tracer sur un axe gradué.
\begin{corrige}[exercice~\ref{t11s-c18-e06}]\label{sol-t11s-c18-e06}
Minimum $2$, $Q_1=x_3=4$, médiane $(5+7)/2=6$, $Q_3=x_8=9$, maximum $12$. Boîte de $4$ à $9$, trait à $6$, moustaches jusqu’à $2$ et $12$.
\end{corrige}
\exo[Des classes]\label{t11s-c18-e07}
Dix trajets durent entre $0$ et $10$ min, vingt entre $10$ et $20$ min. Estimer la moyenne par les centres. Cette valeur est-elle exacte ?
\begin{corrige}[exercice~\ref{t11s-c18-e07}]\label{sol-t11s-c18-e07}
Centres $5$ et $15$; moyenne estimée $(10\times5+20\times15)/30=35/3\simeq11{,}67$ min. Ce n’est pas une moyenne exacte sans les temps individuels; les intervalles doivent être disjoints, par exemple $[0;10[$ et $[10;20]$.
\end{corrige}
\begin{jemeteste}Je vérifie l’effectif total, les rangs, les unités et la distinction entre pourcentage et points de pourcentage.\end{jemeteste}
\begin{notepourlaclasse}Les données sont fictives pour l’apprentissage. La boîte utilise minimum et maximum aux moustaches; ne pas mêler à une convention de détection de valeurs atypiques sans l’annoncer.\end{notepourlaclasse}
